\chapter{अंतःरूपांतरणों का समानयन}\label{ch:b2:reduction}

किसी अंतःरूपांतरण को समझना हो तो वे दिशाएँ खोजिए जिन्हें वह केवल खींचता
है। यह अध्याय उसका यंत्र बनाता है --- \omterm{def:b2:reduction:eigen}{अभिलक्षणिक मान}, अभिलक्षणिक और
\omterm{def:b2:reduction:polyu}{न्यूनतम बहुपद}, केंद्रक अपघटन प्रमेयिका --- और उसके फल भी: विकर्णन तथा
त्रिभुजन की कसौटियाँ, केली--हैमिल्टन, \omterm{thm:b2:reduction:dunford}{डनफ़र्ड अपघटन}, और घातों तथा
घातांकीय फलनों का परिकलन, जिस पर \cref{ch:b2:diffeq} आगे पलेगा। आगे सर्वत्र $E$ एक
परिमित-विमीय $K$-सदिश समष्टि है ($K = \R$ या $\C$) और $u \in
\mathcal{L}(E)$, $n = \dim E$।

\section{अभिलक्षणिक मान और अभिलक्षणिक सदिश}

\begin{definition}\label{def:b2:reduction:eigen}
$\lambda \in K$ $u$ का \emph{अभिलक्षणिक मान}\index{अभिलक्षणिक मान} कहलाता है जब
किसी $x \neq 0$ (जिसे \emph{अभिलक्षणिक सदिश} कहते हैं) के लिए $u(x) = \lambda x$ हो;
\emph{अभिलक्षणिक समष्टि} $E_\lambda(u) =
\ker(u - \lambda\,\mathrm{id})$ है। अभिलक्षणिक मानों का समुच्चय
\emph{वर्णक्रम}\index{वर्णक्रम} $\operatorname{Sp}(u)$ कहलाता है। कोई उपसमष्टि $F$
\emph{स्थायी} कहलाती है जब $u(F) \subseteq F$ हो; अभिलक्षणिक समष्टियाँ स्थायी होती
हैं, और स्थायी उपसमष्टियाँ प्रेरित अंतःरूपांतरण $u|_F$ देती हैं।
\end{definition}

\begin{theorem}[अभिलक्षणिक समष्टियों की स्वतंत्रता]\label{thm:b2:reduction:independence}
परस्पर भिन्न अभिलक्षणिक मानों से जुड़े \omterm{def:b2:reduction:eigen}{अभिलक्षणिक सदिश} स्वतंत्र कुल
बनाते हैं; समतुल्य रूप से, अभिलक्षणिक समष्टियों का योग $E_{\lambda_1} + \dots + E_{\lambda_r}$ (भिन्न $\lambda_i$)
प्रत्यक्ष होता है। विशेष रूप से $u$ के अधिकतम $n$ \omterm{def:b2:reduction:eigen}{अभिलक्षणिक मान} होते
हैं।
\end{theorem}

\begin{proof}
$r$ पर आगमन से। मान लीजिए $x_i \in
E_{\lambda_i}$ सहित $x_1 + \dots + x_r = 0$ है, और कथन $r - 1$ के लिए ज्ञात है।
$u$ लगाइए और उसमें से संबंध का $\lambda_r$ गुना घटाइए:
\[
\sum_{i=1}^{r-1} (\lambda_i - \lambda_r)\, x_i = 0 ,
\]
अतः आगमन से हर $(\lambda_i - \lambda_r)x_i = 0$, अर्थात् $i < r$ के लिए $x_i =
0$, और फिर $x_r = 0$। विमा $n$ की
समष्टि में अशून्य समष्टियों के प्रत्यक्ष योग में अधिकतम $n$ पद होते
हैं।
\end{proof}

\begin{omfigure}
\begin{tikzpicture}[scale=0.9, >={Stealth}]
  \draw[gray!40] (-2.6,-2.2) grid (2.6,2.2);
  \draw[->] (-2.6,0) -- (2.7,0) node[below, font=\small] {$x$};
  \draw[->] (0,-2.2) -- (0,2.3) node[left, font=\small] {$y$};
  \draw[omDef, very thick, ->] (0,0) -- (0.7,0.7)
    node[above, font=\small] {$v_1$};
  \draw[omDef, thick, ->, dashed] (0,0) -- (2.1,2.1)
    node[above right, font=\small] {$Av_1 = 3v_1$};
  \draw[omThm, very thick, ->] (0,0) -- (0.9,-0.9)
    node[below, font=\small] {$v_2$};
  \draw[omThm, thick, ->, dashed] (0,0) -- (0.9,-0.9);
  \draw[omProp, very thick, ->] (0,0) -- (1,0)
    node[below right, font=\small] {$e_1$};
  \draw[omProp, thick, ->, dashed] (0,0) -- (2,1)
    node[right, font=\small] {$Ae_1$};
\end{tikzpicture}

{\small आव्यूह $A = \begin{psmallmatrix}2 & 1\\ 1 &
2\end{psmallmatrix}$ समतल पर क्रिया करते हुए: सामान्य सदिश $e_1$ अपनी रेखा
से हट जाता है, परंतु अभिलक्षणिक दिशाएँ $v_1 =
(1,1)$ और $v_2 = (1,-1)$ केवल खिंचती हैं ---
$3$ से और $1$ से (अतः $Av_2 = v_2$: बिंदुकित प्रतिबिंब $v_2$ से मेल खाता है)।
विकर्णन आधार $(v_1, v_2)$ पर जाना है, जहाँ $A$ $\operatorname{diag}(3, 1)$ बन जाता है।}
\end{omfigure}

\begin{definition}[अभिलक्षणिक बहुपद]\label{def:b2:reduction:charpoly}
$\chi_u(X) = \det(X\,\mathrm{id} - u)$\index{अभिलक्षणिक बहुपद} --- जिसे किसी भी आधार में $\det(XI_n - A)$ के रूप में
परिकलित किया जाता है, और जो घात $n$ का एकिक बहुपद है तथा सादृश्यता के
अंतर्गत अपरिवर्त्य (\cref{thm:b2:linalg:detrules})। $K$ में उसके मूल ठीक \omterm{def:b2:reduction:eigen}{अभिलक्षणिक मान} हैं
($\lambda$ \omterm{def:b2:reduction:eigen}{अभिलक्षणिक मान} $\iff u - \lambda\,\mathrm{id}$ एकैकी नहीं $\iff \chi_u(\lambda) = 0$), और
\[
\chi_u(X) = X^n - (\operatorname{tr} u)\, X^{n-1} + \dots +
(-1)^n \det u .
\]
किसी \omterm{def:b2:reduction:eigen}{अभिलक्षणिक मान} की \emph{बीजीय बहुलता}\index{बीजीय बहुलता} $m_\lambda$
$\chi_u$ के मूल के रूप में उसकी बहुलता है; \emph{ज्यामितीय
बहुलता}\index{ज्यामितीय बहुलता} $\dim E_\lambda$ है, और $1 \leq \dim E_\lambda \leq
m_\lambda$।
\end{definition}

\begin{proof}[कथित तथ्यों की उपपत्ति]
गुणांकों के बारे में दावे: $\det(XI - A)$ को क्रमचय सूत्र से खोलिए; तत्समक क्रमचय
$\prod_i (X - a_{ii})
= X^n - (\sum a_{ii})X^{n-1} + \dots$ का योगदान देता है, और हर दूसरा क्रमचय अधिकतम $n - 2$ विकर्ण स्थान अचर
रखता है, अतः उसका योगदान घात $\leq
n - 2$ का है: इस प्रकार ऊपर के दो गुणांक
कथनानुसार निकलते हैं; और $X = 0$ से अचर पद $\det(-A) = (-1)^n\det A$ मिलता है।

ज्यामितीय $\leq$ बीजीय: मान लीजिए $d = \dim E_\lambda$ और $E_\lambda$ के आधार को $E$ के आधार तक
पूरा कीजिए; $u$ का आव्यूह खंड-रूप में ऊपरी-त्रिभुजाकार है जिसका ऊपरी
बायाँ खंड $\lambda I_d$ है, अतः $\chi_u(X) =
(X - \lambda)^d\, \chi_{\text{(निचला खंड)}}(X)$: अर्थात् $\lambda$ की
बहुलता कम से कम $d$ है।
\end{proof}

\begin{example}[एक ही $\chi$, भिन्न ज्यामिति]\label{ex:b2:reduction:samechi}
आव्यूह
\[
\begin{pmatrix}2 & 0\\ 0 & 2\end{pmatrix}
\qquad\text{तथा}\qquad
\begin{pmatrix}2 & 1\\ 0 & 2\end{pmatrix}
\]
का \omterm{def:b2:reduction:charpoly}{अभिलक्षणिक बहुपद} $(X - 2)^2$, अनुरेख, \omterm{def:b2:linalg:det}{सारणिक} और \omterm{def:b2:reduction:eigen}{वर्णक्रम} एक ही हैं ---
फिर भी वे सदृश नहीं हैं: पहले का $E_2$ विमा $2$ का है (\omterm{def:b2:reduction:charpoly}{ज्यामितीय बहुलता}
$2$), दूसरे का विमा $1$ का। \omterm{def:b2:reduction:charpoly}{अभिलक्षणिक बहुपद} केवल बीजीय बहुलताएँ देखता
है; अभिलक्षणिक समष्टियों की विमाएँ सूक्ष्मतर अपरिवर्त्य हैं, और
\omterm{def:b2:reduction:polyu}{न्यूनतम बहुपद} निर्णय करता है ($X - 2$ बनाम $(X - 2)^2$)। विकर्णनीयता की हर चर्चा
के लिए सार: $\chi$ उम्मीदवारों की सूची बनाता है, पर मत केंद्रक डालते हैं।
\end{example}

\begin{definition}[विकर्णनीय, त्रिभुजनीय]\label{def:b2:reduction:diag}
$u$ \emph{विकर्णनीय}\index{विकर्णनीय} कहलाता है जब $E$ का कोई आधार
अभिलक्षणिक सदिशों से बना हो (आव्यूह के लिए: किसी विकर्ण आव्यूह के
सदृश); और \emph{त्रिभुजनीय}\index{त्रिभुजनीय} तब जब किसी आधार में उसका
आव्यूह ऊपरी त्रिभुजाकार हो।
\end{definition}

\begin{theorem}[विकर्णनीयता की कसौटियाँ]\label{thm:b2:reduction:diagcrit}
निम्नलिखित तुल्य हैं:
\begin{enumerate}
  \item $u$ \omterm{def:b2:reduction:diag}{विकर्णनीय} है;
  \item $E = \bigoplus_{\lambda \in \operatorname{Sp} u} E_\lambda$;
  \item $K$ पर $\chi_u$ पूरी तरह गुणनखंडित होता है और प्रत्येक \omterm{def:b2:reduction:eigen}{अभिलक्षणिक
        मान} के लिए $\dim E_\lambda = m_\lambda$;
  \item (पर्याप्त, आवश्यक नहीं) $K$ में $\chi_u$ के $n$ भिन्न मूल हैं।
\end{enumerate}
\end{theorem}

\begin{proof}
(1 $\iff$ 2): अभिलक्षणिक सदिशों का आधार $E_\lambda$ के आधारों में छँट जाता है,
और विलोमतः प्रत्यक्ष पदों के आधार जोड़ देने से $E$ का आधार बन जाता है
(\cref{thm:b2:reduction:independence} से योग प्रत्यक्ष हो जाता है; और विमाओं की समता से वह सब कुछ बन
जाता है)।

(2 $\iff$ 3): विकर्ण आधार में $\chi_u = \prod (X -
\lambda)^{\dim E_\lambda}$ मेल खाती बहुलताओं के साथ गुणनखंडित हो
जाता है। विलोमतः मान लीजिए $\chi_u$ गुणनखंडित होता है और सर्वत्र $\dim E_\lambda =
m_\lambda$; तब
अभिलक्षणिक समष्टियों का प्रत्यक्ष योग (\cref{thm:b2:reduction:independence} से प्रत्यक्ष) की विमा
\[
\sum_{\lambda}\dim E_\lambda = \sum_{\lambda} m_\lambda =
\deg\chi_u = n ,
\]
है, जिसमें बीच वाली समता इसलिए है कि पूरी तरह गुणनखंडित बहुपद की घात
उसके मूलों की बहुलताओं का योग होती है: अतः योग पूरा $E$ है। ध्यान
दीजिए कि किस परिकल्पना ने क्या किया: गुणनखंडन ने घात भरी, बहुलताओं की
समता ने विमाएँ।

(4 $\Rightarrow$ 1): $n$ भिन्न \omterm{def:b2:reduction:eigen}{अभिलक्षणिक मान} $n$ स्वतंत्र \omterm{def:b2:reduction:eigen}{अभिलक्षणिक सदिश} देते
हैं (\cref{thm:b2:reduction:independence}): अर्थात् एक आधार।
\end{proof}

\begin{method}[विकर्णनीयता का निर्णय]\label{met:b2:reduction:decide}
व्यवहार में इसी क्रम में जाँचिए --- हर चरण काम पूरा कर सकता है। (1)
क्या सरल एवं पूरी तरह गुणनखंडित मूलों वाला कोई विलोपी बहुपद स्वयं सामने
आ रहा है ($u^2 = \mathrm{id}$, $u^2 = u$, $u^k =
\mathrm{id}$)? यदि हाँ: \omterm{def:b2:reduction:diag}{विकर्णनीय}, बिना किसी परिकलन के
(नीचे \cref{cor:b2:reduction:minpolycrit})। (2) $\chi_u$ परिकलित कीजिए; यदि $K$ में उसके $n$ भिन्न मूल
हैं: \omterm{def:b2:reduction:diag}{विकर्णनीय} (\cref{thm:b2:reduction:diagcrit}~(4))। (3) अन्यथा, केवल प्रत्येक बहुगुणित मूल $\lambda$
के लिए $\dim\ker(u -
\lambda\,\mathrm{id})$ की तुलना बहुलता $m_\lambda$ से कीजिए: कहीं भी कमी विकर्णनीयता
समाप्त कर देती है; और सर्वत्र समता उसे सिद्ध कर देती है। सरल मूलों की
अभिलक्षणिक समष्टियाँ कभी परिकलित मत कीजिए (उनकी विमा $1$ होने को
बाध्य है), और केवल निर्णय के लिए कभी त्रिभुजन मत कीजिए।
\end{method}

\begin{example}[विकर्णन काम पर]\label{ex:b2:reduction:diagwork}
$A = I + J = \begin{psmallmatrix}2 & 1 & 1\\ 1 & 2 & 1\\ 1 & 1 &
2\end{psmallmatrix}$, जहाँ $J$ सर्व-एक आव्यूह है:
$\operatorname{Sp}(J) = \{3, 0\}$
से (\cref{ex:b2:linalg:onesmatrix}), $\operatorname{Sp}(A) = \{4,
1\}$, और अभिलक्षणिक समष्टियाँ $\R(1,1,1)$ तथा समतल $\{x + y + z =
0\}$ हैं: विमाएँ
$1 + 2 = 3$, अतः \omterm{def:b2:reduction:diag}{विकर्णनीय} (\cref{thm:b2:reduction:diagcrit}~(2))। बिना किसी आधार-परिवर्तन आव्यूह के घातें:
$\R(1,1,1)$ पर प्रक्षेपक $\Pi = J/3$ लेने पर,
\[
A = 4\,\Pi + 1\cdot(I - \Pi)
\quad\Longrightarrow\quad
A^k = 4^k\,\Pi + (I - \Pi)
= \frac{4^k - 1}{3}\,J + I .
\]
($k = 1$ जाँचिए: $\frac{4-1}3 J + I = A$।) समापन दृष्टि: जब अभिलक्षणिक समष्टियाँ दिखाई देती
हों, तब वर्णक्रमीय \emph{प्रक्षेपक} घातें $PDP^{-1}$ से कहीं तेज़ी से परिकलित
कर देते हैं --- और सूत्र गतिकी भी दिखा देता है: $A^k$ $(1,1,1)$ के अनुदिश
$4^k$ की तरह बढ़ता है और लांबिक समतल पर जहाँ का तहाँ रहता है।
\end{example}

\begin{theorem}[त्रिभुजन]\label{thm:b2:reduction:trigonalization}
$u$ $K$ पर \omterm{def:b2:reduction:diag}{त्रिभुजनीय} है यदि और केवल यदि $\chi_u$ $K$ पर पूरी तरह गुणनखंडित
हो। विशेष रूप से किसी $\C$-सदिश समष्टि का प्रत्येक अंतःरूपांतरण
\omterm{def:b2:reduction:diag}{त्रिभुजनीय} होता है।\index{त्रिभुजन}
\end{theorem}

\begin{proof}
($\Rightarrow$) त्रिभुजाकार आव्यूह का \omterm{def:b2:reduction:charpoly}{अभिलक्षणिक बहुपद} $\prod(X - t_{ii})$ है: अर्थात् पूरी तरह
गुणनखंडित।

($\Leftarrow$) $n$ पर आगमन। चूँकि $\chi_u$ गुणनखंडित होता है, उसका कोई मूल $\lambda$ है:
कोई \omterm{def:b2:reduction:eigen}{अभिलक्षणिक सदिश} $e_1$ चुनिए। $e_1$ से आरंभ होने वाले आधार में आव्यूह
$\begin{pmatrix} \lambda & \ast\\ 0 &
B\end{pmatrix}$ है, और $\chi_u = (X - \lambda)\chi_B$: अतः $\chi_B$ भी गुणनखंडित होता है। $(n-1) \times
(n-1)$ आव्यूह $B$ पर आगमन
परिकल्पना लगाने पर ऐसा प्रतिलोमनीय $Q$ मिलता है कि $Q^{-1}BQ$ ऊपरी
त्रिभुजाकार हो; और $\begin{pmatrix}1 & 0\\
0 & Q\end{pmatrix}$ से पूरे आव्यूह का संयुग्मन करने पर वह
त्रिभुजाकार हो जाता है।
\end{proof}

\begin{example}[हाथ से त्रिभुजन]\label{ex:b2:reduction:trigwork}
$B = \begin{pmatrix}3 & -1\\ 1 & 1\end{pmatrix}$: $\chi_B = X^2 -
4X + 4 = (X - 2)^2$, और $\ker(B - 2I) =
\ker\begin{psmallmatrix}1 & -1\\ 1 & -1\end{psmallmatrix}$ $e_1' = (1, 1)$ से फैली रेखा है: अर्थात् एक \omterm{def:b2:reduction:eigen}{अभिलक्षणिक मान}, और
एक-विमीय \omterm{def:b2:reduction:eigen}{अभिलक्षणिक समष्टि} --- \omterm{def:b2:reduction:diag}{विकर्णनीय} नहीं, पर \omterm{def:b2:reduction:diag}{त्रिभुजनीय} (\cref{thm:b2:reduction:trigonalization})।
आधार को $e_2' = (1, 0)$ से पूरा कीजिए और परिकलित कीजिए:
\[
u(e_1') = (2, 2) = 2e_1',
\qquad
u(e_2') = (3, 1) = 1\cdot e_1' + 2\, e_2' ,
\]
अतः आधार $(e_1', e_2')$ में आव्यूह $T =
\begin{psmallmatrix}2 & 1\\ 0 & 2\end{psmallmatrix}$ है। समापन दृष्टि: $T$ का विकर्ण बाध्य था
(दोनों प्रविष्टियों को दोहरा \omterm{def:b2:reduction:eigen}{अभिलक्षणिक मान} $2$ होना ही था); केवल कोने
वाली प्रविष्टि $e_2'$ के चुनाव पर निर्भर थी, और $e_2'$ का मापन बदलकर उसे
कोई भी अशून्य मान बनाया जा सकता है --- अड़ियल ``$1$'' उस शून्यंभावी
भाग की परछाईं है जिसे डनफ़र्ड अलग करके दिखाएगा।
\end{example}

\section{किसी अंतःरूपांतरण के बहुपद}

\begin{definition}\label{def:b2:reduction:polyu}
$P = \sum a_k X^k \in K[X]$ के लिए $P(u) = \sum a_k u^k \in
\mathcal{L}(E)$ रखिए। प्रतिचित्रण $P \mapsto P(u)$ बीजगणितों की समाकारिता $K[X] \to \mathcal{L}(E)$ है
(\cref{def:b2:structures:algebra}); उसका केंद्रक $\{P : P(u) = 0\}$ $K[X]$ की \omterm{def:b2:structures:ideal}{गुणजावली} है और अशून्य है (कुल $(\mathrm{id}, u, \dots, u^{n^2})$
$n^2$-विमीय $\mathcal{L}(E)$ में परस्पर जुड़ा है), अतः वह किसी अद्वितीय एकिक बहुपद
$\mu_u$ से \omterm{def:b2:structures:generated}{जनित} है: यही \emph{न्यूनतम बहुपद}\index{न्यूनतम बहुपद} है
(\cref{thm:b2:structures:principal})।
\end{definition}

\begin{proposition}\label{prop:b2:reduction:minpoly}
\begin{enumerate}
  \item $P(u) = 0 \iff \mu_u \mid P$; $u$ के \omterm{def:b2:reduction:eigen}{अभिलक्षणिक मान} हर विलोपी बहुपद के मूल होते हैं,
        और $\mu_u$ के मूल \emph{ठीक} \omterm{def:b2:reduction:eigen}{अभिलक्षणिक मान} होते हैं।
  \item यदि $F$ स्थायी है, तो $\mu_{u|_F} \mid \mu_u$।
\end{enumerate}
\end{proposition}

\begin{proof}
(1) विभाज्यता जनक की परिभाषा ही है। यदि $u(x) =
\lambda x$, $x \neq 0$, तो $0 = P(u)(x) = P(\lambda)x$, अतः $P(\lambda) = 0$:
अर्थात् \omterm{def:b2:reduction:eigen}{अभिलक्षणिक मान} विलोपकों के मूल होते हैं, विशेष रूप से $\mu_u$ के।
विलोमतः यदि $\lambda$ कोई मूल है, तो $Q(u) \neq 0$ सहित $\mu_u =
(X - \lambda)Q$ ($\mu_u$ की घात न्यूनतम है):
$Q(u)(y) \neq 0$ वाला $y$ चुनिए; तब $(u - \lambda)(Q(u)(y)) = \mu_u(u)(y)
= 0$ \omterm{def:b2:reduction:eigen}{अभिलक्षणिक सदिश} $Q(u)(y)$ दिखा देता है।

(2) $\mu_u(u|_F) = \mu_u(u)|_F = 0$, और (1) को $u|_F$ पर लगाइए।
\end{proof}

\begin{example}[हाथ से न्यूनतम बहुपद]\label{ex:b2:reduction:minpolyhand}
\omterm{def:b2:reduction:polyu}{न्यूनतम बहुपद} क्रमशः घातें जाँचकर परिकलित किया जाता है। सर्व-एक आव्यूह
$J \in \mathcal{M}_3(\R)$ के लिए: $J \neq \lambda I$ (घात $1$ बाहर हो जाती है), और $J^2 = 3J$, अतः
\[
\mu_J = X^2 - 3X = X(X - 3) :
\]
घात $2$, पूरी तरह गुणनखंडित, सरल मूल --- अतः $J$ \omterm{def:b2:reduction:eigen}{वर्णक्रम} $\{0, 3\}$ के साथ
\omterm{def:b2:reduction:diag}{विकर्णनीय} है (नीचे \cref{cor:b2:reduction:minpolycrit}), और इससे \cref{ex:b2:linalg:onesmatrix} की पुष्टि बिना एक भी \omterm{def:b2:linalg:det}{सारणिक} के
हो जाती है। \cref{ex:b2:reduction:projectorswork} के अदला-बदली आव्यूह $A$ के लिए: $A \neq \pm I$ और $A^2
= I$ से $\mu_A = X^2 - 1$
मिलता है। दोनों स्थितियों में ढर्रा एक ही है: संरचना से किसी कम घात की
सर्वसमिका का अनुमान लगाइए (कोटि एक होने से $J^2 = (\operatorname{tr}J)\,J$ बाध्य हो जाता है; कोई
अंतर्वलन $A^2 = I$ को बाध्य कर देता है), फिर जाँचिए कि कोई उचित भाजक विलोपन
नहीं करता। \omterm{def:b2:reduction:polyu}{न्यूनतम बहुपद} प्रायः \emph{खोजे} जाते हैं, $\chi$ से परिकलित
नहीं किए जाते।
\end{example}

\begin{theorem}[केंद्रक अपघटन प्रमेयिका]\label{thm:b2:reduction:kernels}
यदि $P = P_1 P_2 \cdots P_r$ हो, जहाँ $P_i$ परस्पर सहअभाज्य हैं, तो\index{केंद्रक अपघटन
प्रमेयिका}
\[
\ker P(u) = \ker P_1(u) \oplus \dots \oplus \ker P_r(u),
\]
और पदों पर प्रक्षेप $u$ के बहुपद होते हैं।
\end{theorem}

\begin{proof}
$r = 2$ की स्थिति निपटाकर आगमन कर लेना पर्याप्त है। $K[X]$ में बेज़ू (\cref{thm:b2:structures:principal}):
$U P_1 + V P_2 = 1$, अतः प्रत्येक $x$ के लिए
\[
x = \underbrace{U(u)P_1(u)(x)}_{=:\,x_2}
+ \underbrace{V(u)P_2(u)(x)}_{=:\,x_1}.
\]
यदि $x \in \ker P(u)$: तो $P_2(u)(x_2) = U(u)\,P(u)(x) = 0$ (एक ही $u$ के बहुपद क्रमविनिमेय होते हैं), अतः $x_2 \in \ker P_2(u)$,
और सममित रूप से $x_1
\in \ker P_1(u)$: इस प्रकार योग $\ker P(u)$ को भर देता है; और दोनों पद
$\ker P(u)$ के भीतर बैठते हैं ($P_i \mid P$)। प्रत्यक्षता: $x \in \ker P_1(u)
\cap \ker P_2(u)$ से $x = U(u)P_1(u)x + V(u)P_2(u)x = 0$ मिलता है।
$x_1, x_2$ वाले सूत्र प्रक्षेपों को $V(u)P_2(u)$ और $U(u)P_1(u)$ के रूप में दिखा देते हैं।
\end{proof}

\begin{example}[स्पष्ट प्रक्षेपकों के साथ केंद्रक प्रमेयिका]\label{ex:b2:reduction:projectorswork}
मान लीजिए $A = \begin{psmallmatrix}0 & 1 & 0\\ 1 & 0 & 0\\ 0 & 0 &
1\end{psmallmatrix}$ (पहले दो निर्देशांक आपस में बदल दीजिए)। तब $A^2
= I$: बहुपद
$X^2 - 1 = (X - 1)(X + 1)$ $A$ को विलुप्त करता है, उसके गुणनखंड सहअभाज्य हैं, और बेज़ू स्पष्ट
है:
\[
\frac{1}{2}(X + 1) - \frac12(X - 1) = 1 .
\]
\cref{thm:b2:reduction:kernels} की उपपत्ति के अनुसार $\ker(A - I)$ और $\ker(A + I)$ पर प्रक्षेप $A$ के ये बहुपद हैं:
\[
\pi_+ = \frac{A + I}{2} = \frac12\begin{pmatrix}
1 & 1 & 0\\ 1 & 1 & 0\\ 0 & 0 & 2\end{pmatrix},
\qquad
\pi_- = \frac{I - A}{2} = \frac12\begin{pmatrix}
1 & -1 & 0\\ -1 & 1 & 0\\ 0 & 0 & 0\end{pmatrix}.
\]
जाँच: $\pi_+ + \pi_- = I$, $\pi_+\pi_- = 0$, $\pi_\pm^2 =
\pi_\pm$, और प्रतिबिंब समतल $\{x = y\}$ (सममित सदिश,
\omterm{def:b2:reduction:eigen}{अभिलक्षणिक मान} $1$) तथा रेखा $\R(1, -1, 0)$ (प्रतिसममित, \omterm{def:b2:reduction:eigen}{अभिलक्षणिक मान} $-1$)
हैं। केंद्रक प्रमेयिका केवल अस्तित्व का कथन नहीं है: बेज़ू गुणांक
\emph{ही} प्रक्षेपक सूत्र हैं।
\end{example}

\begin{example}[प्रक्षेपक घातांकीय भी परिकलित कर देते हैं]\label{ex:b2:reduction:expprojectors}
वही अदला-बदली आव्यूह, एक कदम और आगे। चूँकि $A = \pi_+ -
\pi_-$, जहाँ प्रक्षेपक
बीजगणितीय अर्थ में लांबिक हैं ($\pi_+\pi_- = 0$), अतः हर घात $A^k = \pi_+ +
(-1)^k\pi_-$ का पालन करती है,
और घातांकीय श्रेणी प्रक्षेपकों के अनुसार पुनः समूहबद्ध हो जाती है:
\[
\eu^{tA} = \sum_k \frac{t^k}{k!}\bigl(\pi_+ +
(-1)^k\pi_-\bigr)
= \eu^{t}\,\pi_+ + \eu^{-t}\,\pi_- =
\begin{pmatrix}
\cosh t & \sinh t & 0\\
\sinh t & \cosh t & 0\\
0 & 0 & \eu^{t}
\end{pmatrix}.
\]
($t = 0$ जाँचिए: तत्समक; $0$ पर अवकलज: $A$।) अभिलक्षणिक अपघटन किसी आव्यूह
श्रेणी को दो अदिश श्रेणियों में बदल देता है --- यही वह ठीक-ठीक
कार्यविधि है जिसे \cref{ch:b2:diffeq} हर \omterm{def:b2:reduction:diag}{विकर्णनीय} निकाय पर चलाएगा, और यही कारण है कि
सममित युग्मनों का शासन अतिपरवलयिक फलन करते हैं।
\end{example}

\begin{corollary}[न्यूनतम बहुपद से विकर्णनीयता]\label{cor:b2:reduction:minpolycrit}
$u$ \omterm{def:b2:reduction:diag}{विकर्णनीय} है $\iff$ $\mu_u$ $K$ पर \emph{सरल} मूलों के साथ पूरी तरह
गुणनखंडित होता है $\iff$ $u$ का कोई विलोपी बहुपद सरल मूलों के साथ पूरी
तरह गुणनखंडित होता है।
\end{corollary}

\begin{proof}
यदि $P = \prod_{i}(X - \lambda_i)$ (भिन्न $\lambda_i$) सहित $P(u) = 0$, तो प्रमेयिका से $E = \ker P(u) = \bigoplus_i \ker(u -
\lambda_i)$ मिलता है: अर्थात्
अभिलक्षणिक समष्टियों का प्रत्यक्ष योग, अतः $u$ \omterm{def:b2:reduction:diag}{विकर्णनीय} है (\cref{thm:b2:reduction:diagcrit})।
विलोमतः, \omterm{def:b2:reduction:diag}{विकर्णनीय} $u$ $\prod_{\lambda \in \operatorname{Sp}u}(X -
\lambda)$ से मारा जाता है (वह प्रत्येक \omterm{def:b2:reduction:eigen}{अभिलक्षणिक
समष्टि} को मारता है), और वह सरल मूलों के साथ पूरी तरह गुणनखंडित होता है;
और $\mu_u$ उसे विभाजित करता है जबकि उसके मूल वही हैं (\cref{prop:b2:reduction:minpoly}): अतः $\mu_u$ ठीक
वही गुणनफल है।
\end{proof}

\begin{example}\label{ex:b2:reduction:projections}
प्रक्षेप $p^2 = p$ को संतुष्ट करते हैं: अर्थात् $X(X-1)$ से विलुप्त, पूरी तरह
गुणनखंडित सरल मूल --- अतः \omterm{def:b2:reduction:eigen}{वर्णक्रम} $\subseteq \{0, 1\}$ के साथ \omterm{def:b2:reduction:diag}{विकर्णनीय}, और $E =
\ker p \oplus \ker(p - \mathrm{id})$: यही
प्रथम वर्ष का ज्यामितीय विश्लेषण है, जो एक ही पंक्ति में पुनः सिद्ध हो
गया। सममितियाँ ($s^2 = \mathrm{id}$, \omterm{def:b2:linalg:annihilator}{विलोपक} $X^2 - 1$): $\operatorname{char} K
\neq 2$ होने पर \omterm{def:b2:reduction:diag}{विकर्णनीय}, \omterm{def:b2:reduction:eigen}{वर्णक्रम} $\subseteq \{\pm 1\}$।
और $u^3 =
u^2$ तथा $u^2 \neq u$ वाला कोई अंतःरूपांतरण: $X^2(X - 1)$ से विलुप्त, पर
\emph{आवश्यक नहीं} कि \omterm{def:b2:reduction:diag}{विकर्णनीय} हो --- कसौटी इसे पकड़ लेती है (दोहरा
मूल $0$ जाँचना ही पड़ता है: \omterm{def:b2:reduction:diag}{विकर्णनीय} यदि और केवल यदि इसके अतिरिक्त
$\ker u^2 = \ker
u$ भी हो)।
\end{example}

\begin{example}[क्षेत्र निर्णय करता है: $\R^3$ में एक घूर्णन]\label{ex:b2:reduction:rotation3d}
मान लीजिए $R$ $z$-अक्ष के परितः चौथाई घूर्णन है:
\[
R = \begin{pmatrix}
0 & -1 & 0\\
1 & 0 & 0\\
0 & 0 & 1
\end{pmatrix},
\qquad
\chi_R = (X - 1)(X^2 + 1).
\]
$\R$ पर: एकमात्र \omterm{def:b2:reduction:eigen}{अभिलक्षणिक मान} $1$ है, जिसकी \omterm{def:b2:reduction:eigen}{अभिलक्षणिक समष्टि} अक्ष $\R e_3$
है --- अर्थात् अचर सदिशों की एक रेखा, और उससे आगे कोई समानयन नहीं: अतः
$R$ $\mathcal{M}_3(\R)$ में न \omterm{def:b2:reduction:diag}{विकर्णनीय} है न \omterm{def:b2:reduction:diag}{त्रिभुजनीय} ($\chi_R$ गुणनखंडित नहीं होता)।
$\C$ पर: तीन भिन्न \omterm{def:b2:reduction:eigen}{अभिलक्षणिक मान} $1, \iu, -\iu$, अतः $R$ \omterm{def:b2:reduction:diag}{विकर्णनीय} है, जिसके
\omterm{def:b2:reduction:eigen}{अभिलक्षणिक सदिश} $e_3$ और $e_1 \mp \iu e_2$ हैं। ज्यामिति बीजगणित में सुनाई दे रही
थी: समतल में घूर्णनों की कोई वास्तविक अपरिवर्त्य दिशा नहीं होती, और
मापांक $1$ वाले सम्मिश्र \omterm{def:b2:reduction:eigen}{अभिलक्षणिक मान} $\pm\iu$ वह कोण ($\pm\frac\pi2$) सँजोए
रखते हैं जिसे वास्तविक आव्यूह केवल निर्देशांक मिलाकर व्यक्त कर सकता है।
\end{example}

\begin{example}[न्यूनतम बनाम अभिलक्षणिक]\label{ex:b2:reduction:minvschar}
$D = \operatorname{diag}(2, 2, 3)$ के लिए: $\chi_D = (X - 2)^2(X -
3)$, परंतु $\mu_D = (X - 2)(X - 3)$, क्योंकि $(D - 2I)(D - 3I) = 0$ (विहित आधार पर जाँचिए) जबकि
अकेला कोई भी गुणनखंड $D$ को नहीं मारता। स्थानांतरण खंड $N = \begin{psmallmatrix}0 & 1\\ 0 &
0\end{psmallmatrix} \oplus (3)$, अर्थात्
$N' =
\begin{psmallmatrix}0 & 1 & 0\\ 0 & 0 & 0\\ 0 & 0 &
3\end{psmallmatrix}$, के लिए: $\chi_{N'} = X^2(X - 3)$ \emph{और} $\mu_{N'} = X^2(X - 3)$ --- यहाँ दोहरा मूल सचमुच आवश्यक है
क्योंकि $N'$ $\ker$ की ओर \omterm{def:b2:reduction:diag}{विकर्णनीय} नहीं है ($N'e_2 =
e_1 \neq 0$)। अंगूठे का नियम: $\mu$
और $\chi$ के मूल एक ही होते हैं (\cref{prop:b2:reduction:minpoly}); $\mu$ में बहुलता सबसे बड़े
शून्यंभावी खंड का आकार मापती है, और $\chi$ में वाली अभिलक्षणिक उपसमष्टि
की कुल विमा।
\end{example}

\begin{theorem}[केली--हैमिल्टन]\label{thm:b2:reduction:cayleyhamilton}
$\chi_u(u) = 0$; फलतः $\mu_u \mid \chi_u$, और $\deg \mu_u
\leq n$।\index{केली--हैमिल्टन प्रमेय}
\end{theorem}

\begin{proof}
$x \neq 0$ स्थिर कीजिए और $d$ को महत्तम ऐसा लीजिए कि $(x, u(x), \dots,
u^{d-1}(x))$ स्वतंत्र हो; लिखिए
\[
u^d(x) = -a_0 x - a_1 u(x) - \dots - a_{d-1}u^{d-1}(x),
\]
और $P_x = X^d + a_{d-1}X^{d-1} + \dots + a_0$ रखिए, ताकि $P_x(u)(x) =
0$ हो। स्वतंत्र कुल को $E$ के आधार तक पूरा कीजिए:
उस पर $u$ का खंड-रूप $\begin{pmatrix} C & \ast\\ 0 & D\end{pmatrix}$ है, जहाँ $C$ $P_x$ का सहचर आव्यूह है, जिसका
\omterm{def:b2:reduction:charpoly}{अभिलक्षणिक बहुपद} $P_x$ है (पहले स्तंभ के अनुदिश $\det(XI - C)$ खोलिए और $d$ पर
आगमन कीजिए)। अतः $\chi_u = P_x \cdot \chi_D$, और
\[
\chi_u(u)(x) = \chi_D(u)\bigl(P_x(u)(x)\bigr) = 0 .
\]
यह तर्क प्रत्येक $x$ के लिए चलता है: अतः $\chi_u(u) = 0$।
\end{proof}

\begin{example}[केली--हैमिल्टन काम पर]\label{ex:b2:reduction:chwork}
$A = \begin{pmatrix} 1 & 2\\ 3 & 4\end{pmatrix}$: $\chi_A = X^2 -
5X - 2$, अतः $A^2 = 5A + 2I$। $A$ की हर घात $I$ और $A$ के संचय में सिमट जाती है:
\[
A^4 = (5A + 2I)^2 = 25A^2 + 20A + 4I = 145A + 54I =
\begin{pmatrix} 199 & 290\\ 435 & 634\end{pmatrix},
\]
और प्रतिलोम मुफ़्त में मिल जाता है: $A(A - 5I) = 2I$ से
\[
A^{-1} = \tfrac12(A - 5I)
= \begin{pmatrix} -2 & 1\\ 3/2 & -1/2\end{pmatrix}.
\]
समापन दृष्टि: केली--हैमिल्टन पूरे बीजगणित $K[A]$ को $\operatorname{Vect}(I, A, \dots, A^{n-1})$ में संपीड़ित कर
देता है --- $\dim K[A] = \deg\mu_A \leq n$, चाहे आपको कितनी ही बड़ी घातें क्यों न चाहिए।
\end{example}

\begin{remark}[सामान्य भूलें]\label{rem:b2:reduction:pitfalls}
(क) \omterm{def:b2:reduction:eigen}{अभिलक्षणिक मान} जुड़ते नहीं हैं: $\operatorname{Sp}(A + B)$ $\operatorname{Sp}A + \operatorname{Sp}B$ नहीं होता, और \omterm{def:b2:reduction:diag}{विकर्णनीय}
आव्यूहों का योग \omterm{def:b2:reduction:diag}{विकर्णनीय} होना आवश्यक नहीं --- $\begin{psmallmatrix}1 & 1\\ 0 & 0\end{psmallmatrix} +
\begin{psmallmatrix}0 & 0\\ 0 & 1\end{psmallmatrix} =
\begin{psmallmatrix}1 & 1\\ 0 & 1\end{psmallmatrix}$ दो \omterm{def:b2:reduction:diag}{विकर्णनीय}
आव्यूहों का योग है (दोनों के \omterm{def:b2:reduction:eigen}{अभिलक्षणिक मान} भिन्न हैं) और स्वयं
\omterm{def:b2:reduction:diag}{विकर्णनीय} नहीं; केवल \emph{क्रमविनिमेय} कुल ही सधे हुए व्यवहार करते हैं
(\cref{exo:b2:reduction:9})। (ख) ``$\chi_u$ गुणनखंडित होता है'' \emph{क्षेत्र} के बारे में
परिकल्पना है: समतल के घूर्णन का $\chi =
X^2 - 2\cos\theta\,X + 1$ है, जो $\C$ पर गुणनखंडित होता है,
$\R$ पर नहीं --- अतः $\mathcal{M}_2(\C)$ में \omterm{def:b2:reduction:diag}{विकर्णनीय}, $\mathcal{M}_2(\R)$ में \omterm{def:b2:reduction:diag}{त्रिभुजनीय} भी नहीं।
(ग) असमिका सदा ज्यामितीय $\leq$ बीजीय की दिशा में चलती है, कभी उलटी
नहीं; और केवल $\dim E_\lambda \geq
1$ जाँचने से विकर्णनीयता के बारे में कुछ सिद्ध नहीं
होता। (घ) $\mu_u$ $\chi_u$ नहीं है: समता ठीक तब होती है जब प्रत्येक
\omterm{def:b2:reduction:eigen}{अभिलक्षणिक मान} की केवल एक ही खंड-शृंखला हो (जैसे सहचर आव्यूह, इस अध्याय
की सप्ताहांत समस्या); और जहाँ $\mu$ चाहिए वहाँ $\chi$ लगाने से हर घात का
परिकलन फूल जाता है। (ङ) डनफ़र्ड के $d$ और $\nu$ $u$ के बहुपद होते हैं
--- सही गुणधर्मों वाला कोई अपघटन $u = d' + \nu'$ जिसमें $d'\nu' \neq \nu'd'$ हो, डनफ़र्ड
\emph{नहीं} है और कभी अद्वितीय भी नहीं होता।
\end{remark}

\begin{remark}[यह अध्याय कहाँ काम आता है]
समानयन पुस्तक के शेष भाग का श्रमिक है: आव्यूहों की घातें और घातांकीय
\cref{ch:b2:diffeq} के रैखिक अवकल निकाय चलाती हैं; \cref{ch:b2:quadratic} की वर्णक्रमीय प्रमेय विकर्णन का
लांबिक रूप ही है; और जनक फलन (\cref{ch:b2:genfun}) इस अध्याय की सप्ताहांत समस्या वाली
पुनरावृत्ति-अनंतस्पर्शिता को विश्लेषणात्मक ढंग से फिर से निकाल देते
हैं। तृतीय वर्ष के खंड में यही कार्यक्रम अनंत विमा में चलता है: संहत
स्वसंलग्न संकारकों का \omterm{def:b2:reduction:eigen}{वर्णक्रम} सिद्धांत, जहाँ परिमित \omterm{def:b2:reduction:eigen}{वर्णक्रम} का स्थान
अभिलक्षणिक मानों के अनुक्रम ले लेते हैं, और धनात्मक आव्यूहों का
पेरों--फ्रोबेनियस सिद्धांत, जो समझाता है कि गिनती की समस्याओं के प्रबल
\omterm{def:b2:reduction:eigen}{अभिलक्षणिक मान} धनात्मक और सरल \emph{क्यों} होते हैं।
\end{remark}

\section{शून्यंभावी और डनफ़र्ड अपघटन}

\begin{proposition}[शून्यंभावी अंतःरूपांतरण]\label{prop:b2:reduction:nilpotent}
$\chi_u$ के गुणनखंडित होने पर $u$ के लिए निम्नलिखित तुल्य हैं: $u^n =
0$; किसी
$k$ के लिए $u^k = 0$; $\operatorname{Sp}(u) = \{0\}$; $\chi_u
= X^n$; $u$ शून्य विकर्ण के साथ \omterm{def:b2:reduction:diag}{त्रिभुजनीय} है।
किसी \omterm{prop:b2:reduction:nilpotent}{शून्यंभावी अंतःरूपांतरण} के लिए $\mu_u = X^{\text{(शून्यंभाविता
सूचकांक)}}$ होता है, और सूचकांक $\leq n$।\index{शून्यंभावी
अंतःरूपांतरण}
\end{proposition}

\begin{proof}
$u^k = 0$ हर \omterm{def:b2:reduction:eigen}{अभिलक्षणिक मान} को $X^k$ का मूल बना देता है: अतः \omterm{def:b2:reduction:eigen}{वर्णक्रम} $\{0\}$ ($\chi$
के गुणनखंडित होने पर अरिक्त --- $\C$ पर सदा)। तब $\chi_u =
X^n$ (सभी मूल शून्य) और
केली--हैमिल्टन से $u^n = 0$ मिलता है; और त्रिभुजन (\cref{thm:b2:reduction:trigonalization}) विकर्ण पर शून्य रख
देता है (विकर्ण \omterm{def:b2:reduction:eigen}{अभिलक्षणिक मान} वहन करता है)। विलोमतः, मान लीजिए $A$
यथार्थतः ऊपरी त्रिभुजाकार है: $j \leq i$ के लिए $a_{ij} = 0$। आगमन से हम दिखाते हैं कि
\[
(A^k)_{ij} = 0 \qquad \text{ जब भी } j \leq i + k - 1,
\]
अर्थात् हर घात शून्य वाले क्षेत्र को एक विकर्ण ऊपर धकेल देती है। $k = 1$
के लिए यही परिकल्पना है। और चरण के लिए,
\[
(A^{k+1})_{ij} = \sum_{\ell} (A^k)_{i\ell}\,a_{\ell j} ,
\]
जिसमें हर पद लुप्त हो जाता है: या तो $\ell \leq i + k - 1$ (तब पहला गुणनखंड आगमन से $0$
है) अथवा $\ell \geq i + k$, और उस स्थिति में $j \leq i + k \leq \ell$ दूसरे गुणनखंड को मार देता है।
$k = n$ पर शर्त $j \leq i + n - 1$ सभी $i, j \leq n$ के लिए सत्य है: अतः $A^n =
0$। \omterm{def:b2:reduction:polyu}{न्यूनतम बहुपद}
$X^n$ को विभाजित करता है, और विलोपन ही सूचकांक को परिभाषित करता है।
\end{proof}

\begin{theorem}[डनफ़र्ड अपघटन]\label{thm:b2:reduction:dunford}
मान लीजिए $\chi_u$ $K$ पर गुणनखंडित होता है ($K = \C$ के लिए स्वतः)। तब
\[
u = d + \nu, \qquad d \text{ विकर्णनीय}, \quad \nu
\text{ शून्यंभावी}, \quad d\nu = \nu d ,
\]
वाला \emph{अद्वितीय} युग्म $(d, \nu)$ विद्यमान है, और इसके अतिरिक्त $d$ तथा
$\nu$ $u$ के बहुपद होते हैं।\index{डनफ़र्ड अपघटन}
\end{theorem}

\begin{proof}
\emph{अस्तित्व।} $\chi_u = \prod_{i=1}^{r} (X -
\lambda_i)^{m_i}$ लिखिए (भिन्न $\lambda_i$) और $N_i = \ker(u -
\lambda_i)^{m_i}$ रखिए, जिन्हें
\emph{अभिलक्षणिक उपसमष्टि} कहते हैं। केली--हैमिल्टन तथा केंद्रक
प्रमेयिका (\cref{thm:b2:reduction:kernels}) से
\[
E = N_1 \oplus \dots \oplus N_r ,
\]
मिलता है, जिसमें प्रक्षेप $\pi_i$ $u$ के बहुपद हैं; और हर $N_i$ स्थायी है
($u$ के बहुपद $u$ के साथ क्रमविनिमेय होते हैं)। $d = \sum_i \lambda_i
\pi_i$ परिभाषित कीजिए:
यह $u$ का बहुपद है और \omterm{def:b2:reduction:diag}{विकर्णनीय} है (वह $N_i$ पर $\lambda_i$ की तरह क्रिया
करता है, अतः $E$ अपनी अभिलक्षणिक समष्टियों में अपघटित हो जाता है)। तब
$\nu = u -
d$ $u$ का बहुपद है (अतः $d$ के साथ क्रमविनिमेय), और हर $N_i$ पर वह
$u - \lambda_i$ की तरह क्रिया करता है, जहाँ $(u - \lambda_i)^{m_i} = 0$: अर्थात् प्रत्येक पद पर $\nu^{\max m_i} = 0$,
इसलिए $\nu$ शून्यंभावी है।

\emph{अद्वितीयता।} मान लीजिए $u = d' + \nu'$ एक और ऐसा युग्म है। चूँकि $d'$ और $\nu'$
आपस में क्रमविनिमेय हैं, वे $u = d' +
\nu'$ के साथ भी क्रमविनिमेय हैं, अतः $u$ के
हर बहुपद के साथ भी --- विशेष रूप से $d$ और $\nu$ के साथ। तब $d - d'$
\omterm{def:b2:reduction:diag}{विकर्णनीय} है (दो क्रमविनिमेय \omterm{def:b2:reduction:diag}{विकर्णनीय} प्रतिचित्रण एक साथ \omterm{def:b2:reduction:diag}{विकर्णनीय}
होते हैं: \cref{exo:b2:reduction:9}) और वह $\nu' - \nu$ के बराबर है, जो शून्यंभावी है: क्योंकि यदि
$\nu^k = 0$ और $\nu'^{k'} = 0$ हों, तो क्रमविनिमेयता द्विपद प्रसार की अनुमति देती है
\[
(\nu' - \nu)^{k + k' - 1}
= \sum_{j=0}^{k+k'-1}\binom{k + k' - 1}{j}
\,\nu'^{\,j}\,(-\nu)^{k + k' - 1 - j} ,
\]
जिसमें हर पद मर जाता है: या तो $j \geq k'$ (पहला गुणनखंड शून्य) अथवा $k + k' - 1 - j \geq k$
(दूसरा गुणनखंड शून्य), और दोनों में से एक सदा सत्य रहता है। \omterm{def:b2:reduction:diag}{विकर्णनीय}
शून्यंभावी शून्य होता है (उसका \omterm{def:b2:reduction:eigen}{वर्णक्रम} $\{0\}$ है और किसी आधार में वह
विकर्ण है): अतः $d = d'$, $\nu = \nu'$।
\end{proof}

\begin{example}[घातें और घातांकीय]\label{ex:b2:reduction:powers}
$A = \begin{pmatrix} 3 & 1\\ -1 & 1\end{pmatrix}$: $\chi_A = X^2 -
4X + 4 = (X-2)^2$, एकमात्र \omterm{def:b2:reduction:eigen}{अभिलक्षणिक मान} $2$, और विमा $1$ की \omterm{def:b2:reduction:eigen}{अभिलक्षणिक
समष्टि}: अतः \omterm{def:b2:reduction:diag}{विकर्णनीय} नहीं। डनफ़र्ड: $D = 2I$, $N = A - 2I =
\begin{pmatrix} 1 & 1\\ -1 & -1\end{pmatrix}$, $N^2 = 0$। तब
\[
A^k = (2I + N)^k = 2^k I + k\,2^{k-1} N ,
\qquad
\eu^{tA} = \eu^{2t}(I + tN),
\]
जो क्रमशः क्रमविनिमेय द्विपद प्रमेय और घातांकीय श्रेणी (\cref{ch:b2:diffeq}) के
क्रमविनिमेय पदों पर बँट जाने से मिलता है। समानयन आव्यूह गतिकी को अदिश
गतिकी में बदल देता है।
\end{example}

\begin{remark}[इसी खंड के भीतर के परिप्रेक्ष्य]\label{rem:b2:reduction:perspectives}
समानयन एक धुरी है; देखने योग्य चार आरे ये हैं। \cref{ch:b2:nvs} में अनुकूलित
मानदंड ``सभी अभिलक्षणिक मानों का मापांक $< 1$'' को ``किसी संकारक मानदंड
के लिए $< 1$'' में बदल देते हैं, जिससे \omterm{def:b2:reduction:eigen}{वर्णक्रम} घातों और श्रेणियों के
अभिसरण पर शासन करने लगते हैं। \cref{ch:b2:diffeq} में \cref{ex:b2:reduction:powers} की विधि $X' = AX$ का सामान्य हल बन
जाती है: डनफ़र्ड $\eu^{tA}$ को बहुपद-गुणा-घातांकीय खंडों में बाँट देता है, और
स्थायित्व अभिलक्षणिक मानों के वास्तविक भागों से पढ़ लिया जाता है। \cref{ch:b2:quadratic}
में एक अदिश गुणन वह करा देता है जो केवल रैखिक बीजगणित नहीं करा सकता:
सममित आव्यूह \emph{लांबिक रूप से} \omterm{def:b2:reduction:diag}{विकर्णनीय} हो जाते हैं, और उनका
शून्यंभावी भाग बचता ही नहीं। और \cref{ch:b2:genfun} में इस अध्याय की सप्ताहांत समस्या
की प्रबल-अभिलक्षणिक-मान अनंतस्पर्शिता विश्लेषणात्मक रूप में लौटती है,
किसी जनक फलन की सबसे छोटी विचित्रता के रूप में --- एक ही वृद्धि दर की
दो भाषाएँ।
\end{remark}

\section{अभ्यास}

\begin{exercise}[$\star$]\label{exo:b2:reduction:1}
विकर्णन कीजिए (\omterm{def:b2:reduction:eigen}{अभिलक्षणिक मान}, अभिलक्षणिक समष्टियों के आधार,
प्रतिलोमनीय $P$):
\[
A = \begin{pmatrix} 1 & 2\\ 2 & 1 \end{pmatrix},
\qquad
B = \begin{pmatrix} 0 & 1 & 1\\ 1 & 0 & 1\\ 1 & 1 & 0
\end{pmatrix}.
\]
\end{exercise}

\begin{exercise}[$\star$]\label{exo:b2:reduction:2}
दिखाइए कि $C = \begin{pmatrix} 1 & 1\\ 0 & 1\end{pmatrix}$ \omterm{def:b2:reduction:diag}{विकर्णनीय} नहीं है, और दो बार दिखाइए: अभिलक्षणिक
समष्टियों से, और \omterm{def:b2:reduction:polyu}{न्यूनतम बहुपद} से।
\end{exercise}

\begin{exercise}[$\star$]\label{exo:b2:reduction:3}
मान लीजिए $u$ $u^2 - 5u + 6\,\mathrm{id} = 0$ को संतुष्ट करता है। सिद्ध कीजिए कि $u$ \omterm{def:b2:reduction:diag}{विकर्णनीय} है,
संभव \omterm{def:b2:reduction:eigen}{वर्णक्रम} निर्धारित कीजिए, और $u^k$ को $\mathrm{id}$ तथा $u$ के संचय के रूप में
परिकलित कीजिए।
\end{exercise}

\begin{exercise}[$\star\star$]\label{exo:b2:reduction:4}
मान लीजिए $u$ \omterm{def:b2:reduction:diag}{विकर्णनीय} है और $F$ स्थायी उपसमष्टि है। सिद्ध कीजिए कि
$u|_F$ \omterm{def:b2:reduction:diag}{विकर्णनीय} है \emph{(सरल गुणनखंडित मूलों वाले किसी विलोपी बहुपद को
प्रतिबंधित कीजिए)}।
\end{exercise}

\begin{exercise}[$\star\star$]\label{exo:b2:reduction:5}
(फिबोनाच्ची) मान लीजिए $A = \begin{pmatrix} 1 & 1\\ 1 & 0\end{pmatrix}$। $A$ का $\R$ पर विकर्णन कीजिए, और इससे
फिबोनाच्ची अनुक्रम ($F_0 = 0$, $F_1 = 1$, $F_{n+1} = F_n +
F_{n-1}$) के लिए बिने का सूत्र निष्कर्ष रूप
में निकालिए:
\[
F_n = \frac{\varphi^n - \psi^n}{\sqrt 5},
\qquad \varphi = \frac{1 + \sqrt5}{2},\ \psi = \frac{1 -
\sqrt5}{2}.
\]
\end{exercise}

\begin{exercise}[$\star\star$]\label{exo:b2:reduction:6}
मान लीजिए $u \in \mathcal{L}(E)$ है, जहाँ $u^2$ \omterm{def:b2:reduction:diag}{विकर्णनीय} है और $u$ प्रतिलोमनीय ($K = \C$)।
सिद्ध कीजिए कि $u$ \omterm{def:b2:reduction:diag}{विकर्णनीय} है। $u$ के प्रतिलोमनीय न होने पर एक
प्रतिउदाहरण दीजिए।
\end{exercise}

\begin{exercise}[$\star\star$]\label{exo:b2:reduction:7}
निम्नलिखित के लिए \omterm{thm:b2:reduction:dunford}{डनफ़र्ड अपघटन}, $A^k$ और $\eu^{tA}$ परिकलित कीजिए:
\[
A = \begin{pmatrix} 2 & 1 & 0\\ 0 & 2 & 1\\ 0 & 0 & 2
\end{pmatrix}.
\]
\end{exercise}

\begin{exercise}[$\star\star$]\label{exo:b2:reduction:8}
मान लीजिए $A \in \mathcal{M}_n(\C)$ है और किसी $k \geq 1$ के लिए $A^k = I$ है। सिद्ध कीजिए कि $A$
\omterm{def:b2:reduction:diag}{विकर्णनीय} है और उसके \omterm{def:b2:reduction:eigen}{अभिलक्षणिक मान} इकाई के $k$-वें मूल हैं। इससे
निष्कर्ष निकालिए कि इकाई विकर्ण वाले त्रिभुजाकार आव्यूह के सदृश परिमित
कोटि का कोई प्रतिलोमनीय सम्मिश्र आव्यूह तत्समक ही होता है।
\end{exercise}

\begin{exercise}[$\star\star\star$]\label{exo:b2:reduction:9}
(एक साथ विकर्णन) मान लीजिए $u, v$ \omterm{def:b2:reduction:diag}{विकर्णनीय} और परस्पर क्रमविनिमेय हैं।
सिद्ध कीजिए कि वे \emph{एक साथ} \omterm{def:b2:reduction:diag}{विकर्णनीय} हैं: कोई आधार दोनों का
विकर्णन कर देता है। \emph{($u$ की हर \omterm{def:b2:reduction:eigen}{अभिलक्षणिक समष्टि} $v$-स्थायी है;
\cref{exo:b2:reduction:4} का उपयोग करके वहाँ $v$ के प्रतिबंधनों का विकर्णन कीजिए।)}
\end{exercise}

\begin{exercise}[$\star\star\star$]\label{exo:b2:reduction:10}
मान लीजिए $u \in \mathcal{L}(\C^n)$। सिद्ध कीजिए कि $u$ \omterm{def:b2:reduction:diag}{विकर्णनीय} है यदि और केवल यदि
प्रत्येक $u$-स्थायी उपसमष्टि की कोई $u$-स्थायी संपूरक उपसमष्टि हो।
\emph{($\Leftarrow$ के लिए: यह गुणधर्म $F = \sum_\lambda E_\lambda(u)$ पर लगाइए, जो सभी अभिलक्षणिक
समष्टियों का योग है; यदि कोई स्थायी संपूरक $G$ अशून्य होता, तो $u|_G$ का
त्रिभुजन $u$ का कोई \omterm{def:b2:reduction:eigen}{अभिलक्षणिक सदिश} $G$ के भीतर उत्पन्न कर देता --- जो
$G \cap
F = \{0\}$ के विरुद्ध है।)}
\end{exercise}

\begin{exercise}[$\star\star\star$]\label{exo:b2:reduction:11}
(गेलफांड का हल्का रूप --- आगे आने वाले विश्लेषण की $2\times2$ झलक) मान लीजिए
$A \in \mathcal{M}_2(\C)$ है और उसके दोनों अभिलक्षणिक मानों का मापांक $< 1$ है। सिद्ध कीजिए कि
$k \to \infty$ होने पर प्रविष्टि-दर-प्रविष्टि $A^k \to 0$। \emph{(त्रिभुजन कीजिए: $T$
ऊपरी त्रिभुजाकार सहित $A = P(T)P^{-1}$; $T^k$ स्पष्ट रूप से परिकलित कीजिए --- बराबर
और भिन्न अभिलक्षणिक मानों में भेद कीजिए --- और परिबंध लगाइए।)}
\end{exercise}

\begin{exercise}[$\star\star$]\label{exo:b2:reduction:12}
मान लीजिए $\operatorname{rk} u = 1$ ($n
\geq 2$) सहित $u \in \mathcal{L}(\C^n)$ है। दिखाइए कि $\chi_u = X^{n-1}(X - \operatorname{tr} u)$, और यह भी कि $u$
\omterm{def:b2:reduction:diag}{विकर्णनीय} है यदि और केवल यदि $\operatorname{tr} u
\neq 0$। \emph{(\cref{exo:b2:linalg:5} से स्मरण कीजिए कि $u^2 =
(\operatorname{tr} u)\,u$।)}
\end{exercise}

\section{समस्या: रैखिक पुनरावृत्तियाँ और सहचर आव्यूह}

रैखिक पुनरावृत्ति $u_{n+k} = a_{k-1}u_{n+k-1} + \dots + a_0 u_n$ भेस बदली हुई आव्यूह-घात ही है, और समानयन उसे बंद
सूत्रों, वृद्धि दरों तथा त्रुटि आकलनों में बदल देता है। यह सप्ताहांत
समस्या वह शब्दकोश विकसित करती है --- एक ओर सहचर आव्यूह, दूसरी ओर
अनुक्रमों की समष्टि पर स्थानांतरण संकारक --- \emph{रैखिक पुनरावृत्तियों
की मूल प्रमेय} सिद्ध करती है (सामान्य हल \omterm{def:b2:reduction:charpoly}{अभिलक्षणिक बहुपद} के मूलों पर
$\sum_i Q_i(n)\lambda_i^n$ होता है), और उसका लाभ $\sqrt2$ के डायोफैंटीय सन्निकटन पर, पथों और
शब्दों की गिनती पर, तथा युग्मित अनुक्रमों के ऐसे वलय पर खर्च करती है
जिसे केवल एक साथ विकर्णन ही सुलझा सकता है।

\begin{problem}[{सप्ताहांत समस्या --- रैखिक पुनरावृत्तियों की
मूल प्रमेय}]
\label{pb:b2:reduction:1}
$k \geq 1$ स्थिर कीजिए, $a_0
\neq 0$ सहित अदिश $a_0, \dots, a_{k-1} \in \C$, एकिक बहुपद $P = X^k - a_{k-1}X^{k-1} - \dots -
a_1 X - a_0$, और पुनरावृत्ति
\[
(\mathcal R)\colon\quad u_{n+k} = a_{k-1}u_{n+k-1} + \dots +
a_1 u_{n+1} + a_0 u_n \qquad (n \geq 0).
\]
$P$ का \emph{सहचर आव्यूह} यह है:
\[
C =
\begin{pmatrix}
0 & 1 & & \\
 & \ddots & \ddots & \\
 & & 0 & 1\\
a_0 & a_1 & \cdots & a_{k-1}
\end{pmatrix}
\in \mathcal{M}_k(\C).
\]

\medskip
\noindent\textbf{भाग I --- सहचर शब्दकोश।}
\begin{enumerate}
  \item दिखाइए कि कोई अनुक्रम $(u_n)$ $(\mathcal R)$ को संतुष्ट करता है यदि और केवल
        यदि सदिश $v_n = (u_n, u_{n+1}, \dots,
        u_{n+k-1})^{\mathsf T}$ $v_{n+1} = Cv_n$ को संतुष्ट करें, अतः $v_n = C^n v_0$।
  \item सिद्ध कीजिए कि $\chi_C = P$ (पहले स्तंभ के अनुदिश $\det(XI - C)$ खोलिए और $k$ पर
        आगमन कीजिए), और फिर यह भी कि $\mu_C = P$ \emph{($C^{\mathsf T}$ पर जाइए, जिसके लिए
        $e_1$ \omterm{def:b2:structures:generated}{चक्रीय} है, और ध्यान दीजिए कि किसी आव्यूह और उसके \omterm{def:b2:linalg:transpose}{परिवर्त}
        का \omterm{def:b2:reduction:polyu}{न्यूनतम बहुपद} एक ही होता है)}।
  \item दिखाइए कि $P$ के प्रत्येक मूल $\lambda$ के लिए सदिश $(1, \lambda, \dots, \lambda^{k-1})^{\mathsf T}$ $\lambda$ के
        लिए $C$ की \omterm{def:b2:reduction:eigen}{अभिलक्षणिक समष्टि} को फैलाता है; इससे निष्कर्ष
        निकालिए कि $C$ की \emph{हर} \omterm{def:b2:reduction:eigen}{अभिलक्षणिक समष्टि} की विमा $1$ है,
        और $C$ \omterm{def:b2:reduction:diag}{विकर्णनीय} है यदि और केवल यदि $P$ के $k$ भिन्न मूल
        हों।
  \item मान लीजिए $P$ के भिन्न मूल $\lambda_1, \dots,
        \lambda_k$ हैं। दिखाइए कि गुणोत्तर
        अनुक्रम $(\lambda_i^n)_n$ $(\mathcal R)$ के हल-समष्टि का आधार बनाते हैं, अतः हर हल
        अद्वितीय अचरों $c_i$ के साथ $u_n = \sum_i
        c_i\lambda_i^n$ है।
  \item पूरी तरह हल कीजिए: $u_{n+2} = u_{n+1} + 6u_n$, $u_0 = 1$, $u_1 = 8$।
\end{enumerate}

\medskip
\noindent\textbf{भाग II --- स्थानांतरण संकारक और मूल प्रमेय।} मान
लीजिए $\mathcal{S}$ सभी सम्मिश्र अनुक्रमों की $\C$-सदिश समष्टि है और $S \in
\mathcal{L}(\mathcal{S})$
स्थानांतरण है, $S\bigl((u_n)_n\bigr) =
(u_{n+1})_n$।
\begin{enumerate}[resume]
  \item दिखाइए कि $(\mathcal R)$ का हल-समुच्चय $\ker
        P(S)$ है, और उसकी विमा ठीक $k$ है
        \emph{(किसी हल को उसके आरंभिक मानों पर भेजिए)}।
  \item समझाइए कि केंद्रक अपघटन प्रमेयिका (\cref{thm:b2:reduction:kernels}) अनंत-विमीय $\mathcal{S}$ पर $S$
        के लिए बिना किसी परिवर्तन के क्यों लागू होती है, और $P = \prod_{i=1}^{r}(X - \lambda_i)^{m_i}$ के लिए
        $\ker P(S)$ का परिणामी अपघटन लिखिए (भिन्न $\lambda_i$, और सभी अशून्य क्योंकि
        $a_0 \neq 0$)।
  \item $\lambda \neq 0$ और $m \geq 1$ के लिए दिखाइए
        \[
        \ker\,(S - \lambda\,\mathrm{id})^m
        = \bigl\{\,\bigl(Q(n)\,\lambda^n\bigr)_n : Q \in
        \C_{m-1}[X]\,\bigr\},
        \]
        जिसकी विमा $m$ है। \emph{($\Delta Q = Q(X + 1) -
        Q(X)$ के साथ $(S -
        \lambda)\bigl(Q(n)\lambda^n\bigr) =
        \lambda^{n+1}(\Delta Q)(n)$ परिकलित कीजिए, और
        इसका उपयोग कीजिए कि $\Delta$ घात घटा देता है; विमा के लिए आरंभिक
        मानों के द्वारा उसे $m$ से परिबद्ध कीजिए।)}
  \item (रैखिक पुनरावृत्तियों की मूल प्रमेय) निष्कर्ष निकालिए: यदि
        $P = \prod_{i=1}^{r}(X - \lambda_i)^{m_i}$ हो, जहाँ $\lambda_i$ भिन्न और अशून्य हैं, तो $(\mathcal R)$ के हल ठीक वे
        अनुक्रम
        \[
        u_n = \sum_{i=1}^{r} Q_i(n)\,\lambda_i^n,
        \qquad Q_i \in \C_{m_i - 1}[X],
        \]
        हैं जिनमें बहुपद $Q_i$ अद्वितीय रूप से निर्धारित होते हैं।
  \item पूरी तरह हल कीजिए: $u_{n+2} = 4u_{n+1} - 4u_n$, $u_0 =
        1$, $u_1 = 0$, और उत्तर $u_2$ पर जाँचिए।
\end{enumerate}

\medskip
\noindent\textbf{भाग III --- प्रबल मूल और डायोफैंटीय लाभ।}
\begin{enumerate}[resume]
  \item मान लीजिए मूल सरल हैं और $i \geq 2$ के लिए $\abs{\lambda_1} >
        \abs{\lambda_i}$ है, तथा $c_1 \neq 0$ सहित $u_n = \sum_i c_i
        \lambda_i^n$।
        दिखाइए कि $u_n \sim
        c_1\lambda_1^n$ और $u_{n+1}/u_n \to \lambda_1$।
  \item (पेल) $a_{n+1} = a_n + 2b_n$, $b_{n+1} = a_n +
        b_n$, $a_0 = b_0 = 1$ परिभाषित कीजिए। दिखाइए कि $q(a, b) = a^2 - 2b^2$ $q(a_{n+1}, b_{n+1}) = -q(a_n, b_n)$ को
        संतुष्ट करता है, अतः $a_n^2 - 2b_n^2 = (-1)^{n+1}$; और इसका संबंध $M = \begin{psmallmatrix}1 & 2\\ 1 &
        1\end{psmallmatrix}$ के \omterm{def:b2:linalg:det}{सारणिक} से
        बताइए।
  \item इससे त्रुटि आकलन
        \[
        \Bigl|\frac{a_n}{b_n} - \sqrt2\Bigr|
        = \frac{1}{b_n\,(a_n + \sqrt2\,b_n)}
        \leq \frac{1}{2b_n^2},
        \]
        निष्कर्ष रूप में निकालिए और दिखाइए कि वह अनुपात $3 -
        2\sqrt2$ के साथ
        गुणोत्तर ढंग से क्षीण होता है \emph{($M$ के \omterm{def:b2:reduction:eigen}{अभिलक्षणिक मान}
        और $b_n$ की वृद्धि ज्ञात कीजिए)}।
  \item (सामान्य वृद्धि) प्रश्न 9 से सिद्ध कीजिए: (क) यदि हर मूल $\abs{\lambda_i} \leq \rho$
        को संतुष्ट करता है, तो $m = \max_i m_i$ सहित $\abs{u_n} \leq C\,n^{m-1}\rho^n$; (ख) यदि महत्तम मापांक
        का कोई अद्वितीय मूल $\lambda_1$ हो और $Q_1 \neq 0$ हो, तो $u_{n+1}/u_n \to
        \lambda_1$ --- इसे प्रश्न
        10 के हल पर जाँचिए।
\end{enumerate}

\medskip
\noindent\textbf{भाग IV --- पथों और शब्दों की गिनती।} शीर्ष समुच्चय
$\{1, \dots, N\}$ वाले किसी परिमित ग्राफ के लिए \emph{संलग्नता आव्यूह} $A$ में $ij$
के कोर होने पर $A_{ij} = 1$ होता है, अन्यथा $0$।
\begin{enumerate}[resume]
  \item सिद्ध कीजिए कि $(A^n)_{ij}$ $i$ से $j$ तक $n$ लंबाई के पथों की संख्या
        है (अर्थात् $n$ कोरों के अनुक्रम, जिसमें हर पग किसी कोर के
        अनुदिश हो)।
  \item (त्रिभुज) $3$ शीर्षों वाले पूर्ण ग्राफ के लिए $A = J - I$: $J$ के
        \omterm{def:b2:reduction:eigen}{वर्णक्रम} (\cref{ex:b2:linalg:onesmatrix}) का उपयोग करके दिखाइए
        \[
        (A^n)_{ii} = \frac{2^n + 2(-1)^n}{3},
        \qquad
        (A^n)_{ij} = \frac{2^n - (-1)^n}{3} \quad (i \neq j),
        \]
        और $n = 2$ पर पथ गिनकर दोनों जाँचिए।
  \item ($11$ के बिना शब्द) मान लीजिए $w_n$ $n$ लंबाई के उन
        द्विआधारी शब्दों की संख्या है जिनमें दो $1$ लगातार नहीं आते।
        शब्दों को उनके अंतिम अक्षर से संकेतित करके अंतरण आव्यूह पाइए,
        $w_{n+2} = w_{n+1} + w_n$ दिखाइए, $w_n =
        F_{n+2}$ निष्कर्ष रूप में निकालिए (फिबोनाच्ची, \cref{exo:b2:reduction:5}),
        और वृद्धि दर $\lim w_{n+1}/w_n$ दीजिए।
  \item (पथ ग्राफ) पथ ग्राफ $1 - 2 - 3$ के लिए दिखाइए कि $A$ के \omterm{def:b2:reduction:eigen}{अभिलक्षणिक
        मान} $\sqrt2, 0, -\sqrt2$ हैं जिनके \omterm{def:b2:reduction:eigen}{अभिलक्षणिक सदिश} $(1, \pm\sqrt2, 1)$ और $(1, 0, -1)$ हैं, और इससे
        निष्कर्ष निकालिए कि एक सिरे से दूसरे सिरे तक $n$ लंबाई के
        पथों की संख्या $\bigl((\sqrt2)^n + (-\sqrt2)^n\bigr)/4$ है: विषम $n$ के लिए शून्य, और सम $n$ के
        लिए $2^{\,n/2 - 1}$। $n = 4$ पर जाँचिए।
  \item (अनुरेख सूत्र) दिखाइए कि $n$ लंबाई के बंद पथों की कुल संख्या
        (सभी आरंभिक बिंदुओं पर) $\operatorname{tr}(A^n) = \sum_i \lambda_i^n$ है, और इसे त्रिभुज पर सत्यापित
        कीजिए।
\end{enumerate}

\medskip
\noindent\textbf{भाग V --- अनुक्रमों का एक वलय: एक साथ विकर्णन।} $k \geq 3$
स्थिर कीजिए, $\omega = \eu^{2\iu\pi/k}$ लीजिए, और मान लीजिए $W \in \mathcal{M}_k(\C)$ \omterm{def:b2:structures:generated}{चक्रीय} स्थानांतरण है: $W e_i =
e_{i+1}$
(सूचकांक $k$ के सापेक्ष, स्तंभ $0, \dots, k-1$ से अंकित)।
\begin{enumerate}[resume]
  \item दिखाइए कि $W^{\mathsf T}$ $X^k - 1$ का सहचर आव्यूह है, इससे $\chi_W = \mu_W = X^k - 1$ निष्कर्ष रूप में
        निकालिए, और यह भी कि $W$ $k$ सरल अभिलक्षणिक मानों $\omega^j$ तथा
        अभिलक्षणिक सदिशों $f_j = (1, \omega^{-j},
        \omega^{-2j}, \dots, \omega^{-(k-1)j})^{\mathsf T}$ के साथ \omterm{def:b2:reduction:diag}{विकर्णनीय} है।
  \item \emph{\omterm{def:b2:structures:generated}{चक्रीय}} आव्यूह $C = c_0 I + c_1 W + \dots
        + c_{k-1}W^{k-1}$ होता है। दिखाइए कि सभी \omterm{def:b2:structures:generated}{चक्रीय} आव्यूह
        परस्पर क्रमविनिमेय हैं, कि आधार $(f_0, \dots, f_{k-1})$ \emph{उन सबका} एक साथ
        विकर्णन कर देता है, और कि $C$ के \omterm{def:b2:reduction:eigen}{अभिलक्षणिक मान} $\widehat c(\omega^j) = \sum_m
        c_m \omega^{jm}$, $j = 0, \dots, k-1$
        हैं।
  \item इससे $\det C = \prod_{j=0}^{k-1} \widehat
        c(\omega^j)$ निष्कर्ष रूप में निकालिए, और जाँचिए कि $k = 3$ से \cref{exo:b2:linalg:8}
        का गुणनखंडन पुनः प्राप्त हो जाता है।
  \item (माला का औसत) $M = \frac12(W + W^{-1})$ सहित $x^{(n+1)} = Mx^{(n)}$ लीजिए: अर्थात् वलय में सजाई गई
        $k$ संख्याओं में से हर एक को उसके दोनों पड़ोसियों के औसत से
        बदल दिया जाता है। दिखाइए कि $M$ के \omterm{def:b2:reduction:eigen}{अभिलक्षणिक मान} $\cos(2\pi j/k)$ हैं,
        और $f_0$ पर $x^{(0)}$ का गुणांक माध्य $\frac1k\sum_m x^{(0)}_m$ है \emph{($f_j$ के
        निर्देशांकों का योग लीजिए)}।
  \item निष्कर्ष निकालिए: विषम $k$ के लिए $x^{(n)}$ उस अचर सदिश पर अभिसरित
        होता है जिसका मान आरंभिक मानों का माध्य है; और $k = 4$ के लिए वह
        \omterm{def:b2:reduction:eigen}{अभिलक्षणिक मान} दिखाइए जो अनभिसरण के लिए उत्तरदायी है, तथा
        ठीक-ठीक बाधा भी (एक एकांतर-माध्य गुणांक, जिसका लुप्त होना
        आवश्यक है)।
  \item (संश्लेषण) एक-एक वाक्य में: सहचर आव्यूह $(\mathcal R)$ के विश्लेषण को
        समानयन में कैसे बदल देता है; केंद्रक अपघटन प्रमेयिका को कहाँ
        परिमित विमा की आवश्यकता नहीं पड़ी; प्रबल \omterm{def:b2:reduction:eigen}{अभिलक्षणिक मान} वृद्धि
        दरों और डायोफैंटीय त्रुटि पर शासन क्यों करते हैं; संलग्नता
        आव्यूह की घातें पथ क्यों गिनती हैं; और क्रमविनिमेय आव्यूह क्या
        दिला देते हैं। दोनों शिखर नाम लीजिए: रैखिक पुनरावृत्तियों की
        मूल प्रमेय, और --- भाग IV के धनात्मक आव्यूहों के लिए, तृतीय
        वर्ष के खंड में --- पेरों--फ्रोबेनियस प्रमेय।
\end{enumerate}
\end{problem}
